% STAGING ONLY.  This proof is not part of the live manuscript.
% It must not be integrated before exact-byte production Rethlas verification.
\begin{proof}
We separate the part imported from the intersection-line theory from the
new point of the argument.  For a flat projective morphism of relative
dimension \(n\) over a locally Noetherian base, Ducrot constructs a
symmetric multiadditive line
\[
 I_{X/S}(M_0,\ldots,M_n).
\]
Its formation commutes with locally Noetherian base change.  If a section
of one argument cuts out a relative effective Cartier divisor, there is a
canonical restriction isomorphism to the intersection line on that
divisor.  Iterating along a relative regular sequence of length \(n\)
identifies the intersection line with the norm of the remaining argument
on a finite flat scheme.  The restriction isomorphisms commute when two
orders are simultaneously regular.  These are precisely
\cite[Proposition~6.1.1 and Corollary~6.1.2]{Ducrot05Cube}.  We shall use
these established properties, but we reprove the parameter-space and
extension arguments because Ducrot states them for complete sufficiently
positive linear systems.

\medskip
\noindent\emph{The mixed degrees and the line formula.}
Let \(M_1,\ldots,M_n\) be line bundles on \(X\).  Their fibrewise
intersection number can be written as the finite difference
\begin{equation}\label{eq:staging-resultant-euler-difference}
 (M_{1,s}\cdots M_{n,s})_{X_s}
 =
 \sum_{J\subseteq\{1,\ldots,n\}}
 (-1)^{n-|J|}
 \chi\left(
   X_s,\bigotimes_{j\in J}M_{j,s}
 \right).
\end{equation}
Every tensor product in \eqref{eq:staging-resultant-euler-difference} is
flat over \(S\).  Flatness and projectivity make its Hilbert polynomial,
and hence its Euler characteristic, locally constant on \(S\).  This proves
the local constancy of \(d_i\).  If \(n=0\), the empty intersection number
is the length of \(X_s\), which is positive because the fibre is nonempty.
If \(n>0\), the top-dimensional fundamental cycle of a pure fibre is the
sum of its irreducible components with their positive generic lengths.  On
one such component choose an ample bundle \(H\).  For each of the \(n\)
ample bundles \(M_j\), choose \(m_j>0\) such that
\(M_j^{\otimes m_j}\otimes H^{-1}\) is ample.  Ample bundles are nef, mixed
intersections of nef bundles are nonnegative, and expansion of
\[
 \prod_{j=1}^n
 \bigl(c_1(M_j^{\otimes m_j}\otimes H^{-1})+c_1(H)\bigr)
\]
computes
\(\bigl(\prod_jm_j\bigr)(M_1\cdots M_n)\) and contains the strictly
positive term \(c_1(H)^n\).  Here we use the standard intersection-theoretic
facts that mixed intersections of nef Cartier divisors are nonnegative and
that the top self-intersection of an ample Cartier divisor is positive.
It follows that
\((M_1\cdots M_n)>0\) on that component.  Summing with the generic lengths
shows that every \(d_i\) is positive, including on nonreduced fibres.

We use the projection formula for the intersection line.  If
\(\mathscr N\) is a line bundle on the base and \(0\leq i\leq n\), then
\begin{align}
 &I_{Y/P}
 \left(
  A_0,\ldots,A_{i-1},
  A_i\otimes\pi^*\mathscr N,
  A_{i+1},\ldots,A_n
 \right)
 \notag\\
 &\hspace{22mm}\simeq
 I_{Y/P}(A_0,\ldots,A_n)
 \otimes
 \mathscr N^{
  (A_{0,p}\cdots\widehat{A_{i,p}}\cdots A_{n,p})_{Y_p}
 } .
 \label{eq:staging-resultant-intersection-projection}
\end{align}
This follows directly from multiadditivity and the determinant
projection formula; it is also the calculation in
\cite[Lemma~6.3.1]{Ducrot05Cube}.  Apply
\eqref{eq:staging-resultant-intersection-projection} successively to the
base factors \(p_i^*\mathcal O_{P_i}(1)\) in
\eqref{eq:staging-resultant-universal-line}.  Such a base factor is
trivial on a fibre of \(Y/P\), so the exponent at the \(i\)-th step is
exactly \(d_i\), independently of the factors already removed.  Base
change for the intersection line then gives
\[
 I_{Y/P}(g^*L_0,\ldots,g^*L_n)
 \simeq q^*I_{X/S}(L_0,\ldots,L_n).
\]
The last two displays prove
\eqref{eq:staging-resultant-line-formula}.

If \(n>0\), every \(E_i\) has rank at least two.  Indeed, on an
open-and-closed component where \(\operatorname{rk}E_i=1\), the evaluation
quotient is an isomorphism of line bundles.  It would make \(L_i\) trivial
on every fibre, contrary to ampleness on a nonempty positive-dimensional
proper fibre.

\medskip
\noindent\emph{The only replacement for sufficient positivity.}
We record the parameter-space assertion that will be used twice.  Let
\(t\) be a point of a locally Noetherian base \(T\), let
\(W\to T\) be a flat projective scheme obtained, after base change from
\(X\), by intersecting a relative regular subsequence of the universal
sections, and suppose that
the fibres of \(W/T\) have positive dimension.  Fix an unused index \(i\).
For each \(x\in\operatorname{Ass}(W_t)\), evaluation induces a nonzero
\(\kappa(t)\)-linear map
\begin{equation}\label{eq:staging-resultant-associated-evaluation}
 E_{i,t}\longrightarrow
 L_i\otimes_{\mathcal O_X}\kappa(x).
\end{equation}
Indeed, the surjection \eqref{eq:staging-resultant-evaluation} remains
surjective after restriction to \(W_t\) and scalar extension to
\(\kappa(x)\); hence at least one element of \(E_{i,t}\) has nonzero value
at \(x\).  The kernel of
\eqref{eq:staging-resultant-associated-evaluation} is therefore a proper
linear subspace of \(E_{i,t}\).

There are only finitely many associated points.  The fibrewise criterion
for a relative effective Cartier divisor says that a section is regular on
\(W_t\) exactly when it avoids all of them.  Consequently the parameters
for which the next universal section is a zero divisor are contained in a
finite union of proper projective linear subspaces of \(P_{i,t}\).
There is no additional open set of nowhere-vanishing sections to remove.
If a section of \(L_i|_{W_t}\) were nowhere vanishing, it would trivialize
that line bundle.  But \(L_i|_{W_t}\) is ample, whereas the trivial bundle
on a nonempty proper scheme with a positive-dimensional component is not
ample.  Thus, whenever the preceding intersection has positive dimension,
a section which is regular also has a nonempty zero scheme.  We have proved
the analogue, for the quotient system \(E_i\), of
\cite[Lemmas~2.4.1 and~2.4.3]{Ducrot05Cube}.

The same fibrewise Cartier criterion shows that relative regularity is an
open condition on the parameter space.  Moreover, every irreducible
component of a nonempty effective Cartier divisor has codimension one in an
irreducible component of the preceding pure-dimensional fibre.  The divisor
is therefore pure of dimension one less, and induction shows that every
nonempty intersection cut by a relative regular subsequence remains pure of
the expected dimension.

This is also the point at which ampleness is essential to the present
packet.  Basepoint-freeness alone only proves that the kernels in
\eqref{eq:staging-resultant-associated-evaluation} are proper; it does not
exclude a nowhere-vanishing section on a positive-dimensional preceding
intersection.

\medskip
\noindent\emph{Depth of the two exceptional loci.}
For a permutation \(\tau\), let \(U_\tau\) be as in the statement, and put
\[
 U:=\bigcup_\tau U_\tau,
 \qquad
 V:=\bigcap_\tau U_\tau,
 \qquad
 Z:=P\setminus U,
 \qquad
 Z':=P\setminus V .
\]
We claim that
\begin{equation}\label{eq:staging-resultant-two-depth-bounds}
 \begin{aligned}
 z\in Z
 &\Longrightarrow
 \operatorname{depth}\mathcal O_{P_{q(z)},z}\geq2,\\
 z\in Z'
 &\Longrightarrow
 \operatorname{depth}\mathcal O_{P_{q(z)},z}\geq1.
 \end{aligned}
\end{equation}

Fix \(z\in Z\), put \(s=q(z)\), and choose among the universal sections an
ordered relative regular sequence at \(z\) of maximal length \(k\).  Since
\(z\notin U\), one has \(k<n\), so at least two indices remain unused.
Let \(B\) be the product over \(\kappa(s)\) of the \(k\) projective factors
that occur in this sequence, let \(\rho\colon P_s\to B\) be the projection,
and put \(t=\rho(z)\).  Faithfully flat descent along
\(\kappa(t)\to\kappa(z)\) shows that the chosen sections cut out on
\(X_s\times_{\kappa(s)}\kappa(t)\) a nonempty regular intersection \(W_t\)
of pure dimension \(n-k>0\).

For an unused index \(i\), let \(C_i\subset P_{i,t}\) be the bad parameter
locus: its points are those for which the next section does not extend the
chosen sequence to a regular sequence on \(W_t\).  It is closed, and the
parameter-space assertion above places it in a finite union of proper
projective linear subspaces; in particular every irreducible component of
\(C_i\) has codimension at least one in \(P_{i,t}\).  Maximality, together
with preservation and reflection of regularity under the residue-field extension
\(\kappa(t)\to\kappa(z)\), places the image of \(z\) in \(C_i\) for every
unused \(i\).

The fibre
\[
 F_t:=\rho^{-1}(t)
 \simeq\prod_{i\,\mathrm{unused}}{}_{\kappa(t)}P_{i,t}
\]
is regular; here and below \(z\) also denotes its induced point of \(F_t\).
For two distinct unused indices \(i,j\), the inverse images of
\(C_i\) and \(C_j\) in \(F_t\) involve independent projective factors;
their intersection has codimension at least two and contains \(z\).  Hence
\(\operatorname{depth}\mathcal O_{F_t,z}\geq2\).  The projection \(\rho\)
is smooth, so the flat local depth inequality gives
\[
 \operatorname{depth}\mathcal O_{P_s,z}
 \geq \operatorname{depth}\mathcal O_{F_t,z}
 \geq2.
\]

If \(z\in Z'\), choose a permutation for which the indicated sequence is
not regular and stop immediately before its first failure.  The preceding
intersection still has positive dimension.  Projecting to the factors that
precede the failure and then passing to the fibre over the image of \(z\),
the corresponding bad locus \(C_i\) has codimension at least one in the
next projective factor and contains \(z\).  The same fibre-depth and flat
local depth argument gives
\(\operatorname{depth}\mathcal O_{P_s,z}\geq1\).  This proves
\eqref{eq:staging-resultant-two-depth-bounds}.  It is the proof of
\cite[Lemma~6.3.2]{Ducrot05Cube}, with the complete space of sections
replaced by \eqref{eq:staging-resultant-associated-evaluation}.

For clarity, the associated-prime avoidance used here is not by itself an
iteration certificate.  After every cut, the associated points are those
of the new quotient, and the argument above recomputes them.  The positive
dimension of the quotient, together with ampleness, rules out an empty
next cut.  This is exactly the quotient-side hypothesis that would be
missing if one tried to infer a whole regular sequence from avoidance on
\(X_s\) alone.

\medskip
\noindent\emph{Local norm sections, gluing, and extension.}
On \(U_\tau\), successive restriction in the intersection-line functor
gives
\begin{equation}\label{eq:staging-resultant-successive-restriction}
 \mathscr R|_{U_\tau}
 \simeq
 N_{W_\tau/U_\tau}
 \left(\widetilde L_{\tau(n)}|_{W_\tau}\right),
 \qquad
 W_\tau:=
 \bigcap_{j=0}^{n-1}V(\sigma_{\tau(j)}).
\end{equation}
The scheme \(W_\tau\to U_\tau\) is finite and flat.  Define \(r_\tau\) to
be the norm of \(\sigma_{\tau(n)}|_{W_\tau}\) under
\eqref{eq:staging-resultant-successive-restriction}.

On \(V\), every ordering is regular.  Compatibility of two successive
restriction maps and symmetry of the intersection line show that
\(r_\tau=r_\upsilon\) there for all \(\tau,\upsilon\); this is the
commutative restriction calculation of
\cite[Lemma~6.3.3]{Ducrot05Cube}.  We explain why equality extends to the
whole overlap.  The morphism \(P\to S\) is smooth.  The flat local depth
inequality and the second line of
\eqref{eq:staging-resultant-two-depth-bounds} give
\[
 \operatorname{depth}\mathcal O_{P,z}\geq1
 \qquad(z\in Z').
\]
After trivializing \(\mathscr R\) on an open subset of
\(U_\tau\cap U_\upsilon\), the difference \(r_\tau-r_\upsilon\) is a
regular function supported in \(Z'\).  Depth at least one says that
\(\mathcal H^0_{Z'}(\mathcal O_P)=0\), so the difference is zero.  The \(r_\tau\)
therefore glue to a section \(r_U\in H^0(U,\mathscr R)\).

The first line of \eqref{eq:staging-resultant-two-depth-bounds} and the same
flat local depth inequality give
\[
 \operatorname{depth}\mathcal O_{P,z}\geq2
 \qquad(z\in Z).
\]
Locally trivialize \(\mathscr R\).  The local-cohomology sequence for
\(j\colon U\hookrightarrow P\) begins
\[
 0\longrightarrow \mathcal H^0_Z(\mathcal O_P)
 \longrightarrow\mathcal O_P
 \longrightarrow j_*\mathcal O_U
 \longrightarrow \mathcal H^1_Z(\mathcal O_P).
\]
The pointwise depth bound, applied on each local open subset meeting \(Z\),
makes both local-cohomology sheaves vanish.
Hence every local coefficient of \(r_U\) extends uniquely across \(Z\),
and the extensions agree on overlaps.  This produces the unique global
section in \eqref{eq:staging-resultant-section}.  Its construction gives
the asserted local norm characterization.  No normality, reducedness, or
\(S_2\) hypothesis on \(S\) has been used.

\medskip
\noindent\emph{Incidence at fibre-local depth at most one.}
The intersection \(\bigcap_iV(\sigma_i)\) is closed in \(Y\), and
\(\pi\) is projective, so its image \(\mathcal Z\) is closed.  Let
\(p\in P\) satisfy
\(\operatorname{depth}\mathcal O_{P_{q(p)},p}\leq1\).
The first line of \eqref{eq:staging-resultant-two-depth-bounds} implies
\(p\in U\).  Choose \(\tau\) with \(p\in U_\tau\).  On the fibre at \(p\),
the norm in \eqref{eq:staging-resultant-successive-restriction} is the
determinant of multiplication by a representative of
\(\sigma_{\tau(n)}\) in a finite-dimensional
\(\kappa(p)\)-algebra.  That determinant is zero exactly when the
representative is not a unit, equivalently when it belongs to a maximal
ideal of the finite algebra.  This is equivalent to the last section
vanishing at a point of the finite intersection cut out by the first
\(n\) sections.  We obtain
\eqref{eq:staging-resultant-incidence-detection}.  When \(n=0\), the empty
regular sequence gives \(U=P\).  In this case \(Y\to P\) is projective with
zero-dimensional fibres, hence proper and quasi-finite, therefore finite;
it is flat by hypothesis.  The same finite-flat norm argument consequently
applies at every point, without a depth restriction.

\medskip
\noindent\emph{Symmetry and base change.}
Permuting the factors carries the universal sections to the permuted
universal sections.  On every \(U_\tau\), compatibility of successive
restriction with the symmetry constraints of \(I_{Y/P}\) carries the norm
description to the corresponding permuted norm description.  Uniqueness
of extension across \(P\setminus U\) proves symmetry of the global
section.

Let \(a\colon T\to S\) be a morphism with \(T\) locally Noetherian.  Projective
bundles, tautological sublines, universal sections, intersection lines,
successive restriction, and finite-flat norms all commute with this base
change.  The open \(U_\tau\) also commutes with base change.  Indeed,
relative regularity is checked on fibres; after passing from a point of
\(P\) to a point above it, the relevant residue-field extension is
faithfully flat, and it preserves and reflects non-zero-divisors,
successive quotients, and nonemptiness.  Thus the pulled-back section and
the section constructed over \(T\) have the same norm description on the
base change of \(U\).  The depth-two uniqueness argument over \(T\) makes
them equal globally.  Formula
\eqref{eq:staging-resultant-euler-difference} proves
\(d_{i,T}=a^*d_i\), and the construction of
\eqref{eq:staging-resultant-line-formula} proves compatibility of that
isomorphism.

\medskip
\noindent\emph{The rank-one endpoint.}
The assertion for \(n>0\) was proved before the parameter-space argument.
If \(n=0\) and \(\operatorname{rk}E_0=1\), then \(P_0=S\) and
\(\mathcal O_{P_0}(-1)=E_0\).  The evaluation is an isomorphism, so
\[
 \widetilde L_0
 =
 g^*L_0\otimes\pi^*E_0^\vee
 \simeq\mathcal O_Y,
 \qquad
 \sigma_0=1.
\]
In relative dimension zero the intersection line is the finite-flat norm.
It follows that \(\mathscr R\simeq\mathcal O_P\) and that the resultant is
\(1\).  Under \eqref{eq:staging-resultant-line-formula}, the factor
\(I_{X/S}(L_0)=N_{X/S}(f^*E_0)\) is
\(E_0^{\otimes d_0}\), while
\(\mathcal O_{P_0}(d_0)=(E_0^\vee)^{\otimes d_0}\); the two factors cancel.
This also checks the sign of the exponent in
\eqref{eq:staging-resultant-line-formula}.
\end{proof}

\begin{example}\label{ex:staging-resultant-codimension-two-firewall}
Let \(k\) be a field, let
\(X=\mathbf P^1_k\times\mathbf P^1_k\), and take three copies of the
basepoint-free system
\[
 L_i=\operatorname{pr}_1^*\mathcal O_{\mathbf P^1}(1),
 \qquad
 E_i=H^0(\mathbf P^1,\mathcal O_{\mathbf P^1}(1)).
\]
The parameter space is \(P=(\mathbf P^1)^3\).  Each nonzero member vanishes
on \(\{a_i\}\times\mathbf P^1\) for a unique \(a_i\in\mathbf P^1\), so the
three sections have a common zero exactly on the small diagonal
\[
 \{a_0=a_1=a_2\}\subset(\mathbf P^1)^3.
\]
This locus has codimension two.  Moreover every mixed degree is
\[
 \bigl(\operatorname{pr}_1^*c_1(\mathcal O(1))\bigr)^2=0,
\]
so every prescribed multidegree of the cycle resultant is zero.  Hence the
normalized cycle resultant is constant; evaluation at any tuple having no
common zero normalizes that constant to \(1\).  At the generic point of
the small diagonal the local depth of \(P\) is two.  Thus a statement
extending the theorem to semiample systems while deleting
\eqref{eq:staging-resultant-depth-firewall} would be false: without
ampleness, a divisorial resultant need not detect every
higher-codimensional common-zero locus.  The example does not rule out a
stronger all-point detection statement under the ample hypotheses of the
theorem; that separate strengthening is not claimed here.
\end{example}
